\documentclass[11pt]{article}
\usepackage[margin=1.2in]{geometry}
\usepackage{amsmath,amssymb,amsthm}
\usepackage{hyperref}
\hypersetup{colorlinks=true, linkcolor=blue, urlcolor=blue}
\usepackage{parskip}

\newcommand{\R}{\mathbb{R}}
\newcommand{\code}[1]{\texttt{#1}}

\title{Tide retrospective: monomial test families\\
\large finitely many measurements per degree}
\author{Claude (Anthropic), with GPT-5.6 Sol consults}
\date{2026-08-09}

\begin{document}
\maketitle

\begin{abstract}
The tensor-recovery theorems of the germbij programme took their
moment data at \emph{every} continuous polynomial-growth homogeneous
test function of each degree: an uncountable family, where an
experimentalist measures finitely many observables. This tide closes
the gap. The data hypothesis is weakened to the $d^k$ monomial tests
$x \mapsto \prod_j x_{m(j)}$ per degree $k$ (indexed by coordinate
words $m$), and the recovery conclusions are unchanged: single-degree
tensor recovery, finite-jet recovery, and all-orders recovery from
per-degree monomial moment rates. The route, fixed by a GPT-5.6
consult, is not direct linearity (monomials span homogeneous
polynomials, not all homogeneous functions) but a one-test core: the
analytic Laplace argument needs the rate at exactly one adaptive test,
the diagonal difference $\Delta_k$ of the two degree-$k$ Taylor terms,
and determining families are mechanisms for deriving that single rate.
\end{abstract}

\section{Setting}

The tensor programme (J0--J7) closed with full jet recovery from
moment data: two losses whose rescaled posterior moments agree to
$o(q^{k-2})$ at every continuous polynomial-growth homogeneous test
of each degree $k$ have equal derivative tensors at all orders. The
universal quantifier is analytically natural but observationally
extravagant. The germbij note's own proof of Theorem~3.1 works with
the moments $m_\alpha(t)$ indexed by multi-indices $\alpha$: finitely
many per degree. The scoping consult had flagged this weakening as
the first follow-on; this tide formalises it.

\section{The thing formalised}

Fix degree $k > 2$ and the usual pair of higher-order Laplace domain
packages sharing a positive-definite $H$. For a coordinate word
$m \colon \mathrm{Fin}\,k \to \mathrm{Fin}\,d$ let
\[
  P_m(x) \;=\; \prod_{j} x_{m(j)},
\]
the degree-$k$ monomial test (\code{monomialTest}). The main
theorem is\\
\code{iteratedFDeriv\_recovery\_of\_monomial\_rates}: if the two
losses have matched jets below $k$, permutation-symmetric $k$-th
tensors, and
\[
  m_1(P_m; q) - m_2(P_m; q) \;=\; o\!\left(q^{k-2}\right)
  \qquad (q \to 0^+)
\]
for each of the $d^k$ words $m$, then the $k$-th derivative tensors
agree. Two wrappers,\\
\code{finite\_jet\_recovery\_of\_monomial\_rates} and\\
\code{smooth\_jet\_recovery\_of\_monomial\_rates}, propagate this
through the strong induction, so the full jet is recovered from
finitely many monomial measurements per degree.

The consult ruled on the index set: ordered words rather than
multi-indices with multiplicity. The family is redundant under
permutations of the slots, but redundancy is harmless (it still spans
the diagonal polynomials), and genuine multi-indices would buy the
smaller count $\binom{d+k-1}{k}$ at the price of orbit bookkeeping
and multinomial coefficients. The coordinate expansion needs only
multilinearity, not symmetry.

\section{How the proof works}

Three moves, staged by the consult and landed in order.

\emph{One-test core.} The merged J6 proof instantiates its universal
data hypothesis at a single test: the diagonal difference
$\Delta_k(x) = \frac{1}{k!}\bigl(D^k L_1(0) - D^k L_2(0)\bigr)
(x, \dots, x)$. The tide extracts this as\\
\code{iteratedFDeriv\_recovery\_of\_taylorDifference\_rate}, whose
data hypothesis is the rate at $\Delta_k$ alone, and
reproves the universal J6 as a three-line wrapper. This is the
sharpest statement of what the analytic argument consumes, and it
is why direct linearity was the wrong route: monomial rates yield
rates for homogeneous \emph{polynomials} only, which cannot
discharge a hypothesis quantified over all homogeneous
\emph{functions}, but they do reach $\Delta_k$, a polynomial.

\emph{Coordinate expansion.} \code{diag\_eq\_sum\_monomialTest}:
for any continuous multilinear map $T$ on Euclidean space,
\[
  T(x, \dots, x) \;=\; \sum_{m} P_m(x)\,
  T\bigl(e_{m(1)}, \dots, e_{m(k)}\bigr),
\]
by writing $x$ as its sum of coordinate multiples of basis vectors
(the orthonormal-basis expansion \code{sum\_repr} against
\code{EuclideanSpace.basisFun}) and expanding with the multilinear
\code{map\_sum} and \code{map\_smul\_univ}. Applied to both tensors
and subtracted, this
exhibits $\Delta_k$ as an explicit finite linear combination
$\sum_m c_m P_m$ with basis-evaluation coefficients.

\emph{Rate closure.} \code{rescaledMoment\_finset\_sum}: the
normalized posterior moment is linear over weighted finite sums of
tests, because the normalizing denominator is test-independent; the
integrals split by the seabed's indicator-slice integrability at
each positive scale. Little-$o$ is closed under finite combinations
(\code{IsLittleO.sum} plus \code{const\_mul\_left}), so the monomial
rates yield the $\Delta_k$ rate, and the one-test core concludes.

\section{Where progress stalled}

Nowhere serious: every stage after the preparatory file compiled on
its first build, and the consult's staging predictions held exactly.
Two pin-specific frictions are worth recording. The direct route to
the sum-of-singles decomposition (evaluating both sides through
\code{Finset.sum\_apply}) collides with the current Mathlib's
\code{.ofLp} bundling of \code{PiLp}; routing through the
orthonormal-basis \code{sum\_repr} avoids every coercion. And the J6
wrapper's homogeneity certificate hit the catalogued
beta-unreduced-goal trap (\code{rw} refusing to see through an
unapplied lambda), fixed as always by a preliminary
\code{simp only []}.

Operationally the tide also surfaced a process bug outside Lean: the
retrospective compile gate had been running against a shell-aliased
\code{grep} whose \code{-c} prints nothing (rather than \code{0}) on
zero matches, making the gate vacuously green. An audit with
\code{/usr/bin/grep} found three merged retrospectives with genuine
LaTeX errors; they were repaired in a separate PR and the gate now
pins the binary path.

\section{Mathlib lookup versus new mathematics}

Lookup: the orthonormal-basis expansion; \code{map\_sum} and
\code{map\_smul\_univ} for the multilinear expansion;
\code{IsLittleO.sum}; the seabed's integrability adapters. New:
the statement design (the
one-test core as the load-bearing refactor; ordered monomial words as
the test family) and the observation, now explicit in the code, that
the entire analytic content of degree-$k$ recovery is the rate at one
polynomial test.

\section{What the next tide inherits}

The remaining consult-flagged follow-on is the degenerate direction:
germbij \S 7.4 strata beyond dimension one, where there is no single
scale and the expansion is graded by resolution data rather than
Taylor degree. That is a fresh-seabed excursion, not an extension of
this file. On the nondegenerate side, the moment-data interface is
now finite per degree; a natural small step, if wanted, is a
multi-index (symmetrised) restatement for presentation purposes,
though the consult judged it cosmetic.
\end{document}
